\section{Team Collaboration}
\label{sec:collaboration}

\paragraph{Roles.} Table~\ref{tab:collab-roles} lists each member's contributions.

\begin{table}[H]\centering\small
\captionsetup{justification=raggedright,singlelinecheck=false}

\begin{tabular}{@{}L{.20\linewidth}L{.14\linewidth}L{.60\linewidth}@{}}
\toprule
\textbf{Name} & \textbf{Student ID} & \textbf{Contributions} \\
\midrule
Kiatisak Markmeeshap & 67070501005 & Tasks 1 to 5: designed the overall approach and the model pipelines, researched techniques and planned the experiments that the team ran, did the EDA, set up the first version of the report, and assembled the final submission file. \\
Thanaboon Tikaew & 67070501021 & Task 4 and its EDA: looked for structure in the screenshots and found the near-duplicate scenes inside the training set that make random folds over-optimistic, along with the 4:1 class imbalance and the overlap of brightness and saturation between titles; trained and tested a DINOv2-L fine-tune as an extra ensemble member; wrote the Task 4 pipeline-progression and evaluation sections. \\
Chanon Lhumsa-ard & 67070501059 & Tasks 3 and 5 and their EDA: in Task 3, showed that historical spending supplies the offset for 3,112 of the 3,125 payers without a previous-month value and documented the preprocessing; in Task 5, explained the per-period check of the 3-metric box, traced the low-stability recovery from 1,257 accounts to 53 flags, and re-checked the detector against the final 2,104-account file. Also ran Task 4 GPU jobs (ConvNeXt-B and ConvNeXt-L DINOv3, WiSE-FT, frozen-feature kernels) and wrote the first Task 1 report, which set the structure of every chapter. \\
Siriwan Yindeephot & 67070501075 & Tasks 1 and 2 and their EDA: in Task 1, linked the missing-feature analysis (OOF AUC 0.922 with all key features, 0.831 with two or more missing) to how the model should be monitored and deployed, and wrote the business-impact section; in Task 2, clarified how discrete features enter the Naive Bayes and the discriminative models and compared the tree baseline with Naive Bayes. Also trained the EVA02-S and EVA02-B models for the Task 4 ensemble. \\
\bottomrule
\end{tabular}
\caption{Team roles.}\label{tab:collab-roles}
\end{table}

\paragraph{What the logs show.} Several teammates ran long GPU jobs on their own Kaggle and Colab accounts so that we could use more than one account's quota: for example the EVA02-B fine-tune, the third fold of the ConvNeXt-B run, a DINOv2-L fold, frozen-feature extraction kernels, and a WiSE-FT experiment. A teammate wrote the first draft of the Task 1 report (pull request \#1 of the submission repository), which set the structure that this report follows for all five tasks. Questions about the rules were routed to the TA through one teammate, and the answers (15 and 16 November) shaped the final submission.

\paragraph{Key takeaways from teamwork.}
\begin{itemize}[leftmargin=*,itemsep=2pt]
\item \textbf{Share scarce compute across accounts.} GPU quota is per account, and some notebooks need credentials that only a UI session has (for example gated DINOv3 weights, which a command-line push cannot read). Agreeing who runs which job, and handing results over as files in a fixed folder layout, avoided idle sessions and duplicated runs.
\item \textbf{Treat data re-releases and row order as a risk.} On 12 November the organisers re-released the Task 2, 4 and 5 data. Our Task 4 predictions had been written in the old test order, and the first submissions on the new file scored at chance until we re-aligned every prediction by file name. We now key every prediction file by an id or file name and check it in \texttt{combine\_submission.py}, never by row position alone.
\item \textbf{Route rule questions through one person to the TA.} One channel kept the questions precise and the answers consistent. The rulings of 15 and 16 November changed the design (no test-based prior, no test-rank threshold), and the team could check every candidate file against the same text.
\item \textbf{Change one thing per leaderboard submission.} Each probe changed exactly one task column and had its expected effect written down first, so a score difference could be attributed to that change. The same discipline showed which late changes were below the leaderboard's resolution.
\item \textbf{Keep one canonical log.} A single running file with every submission, its change and its score let any teammate see what had been tried and prevented repeating failed experiments; the failed experiments in each task section are part of the result.
\end{itemize}
